\documentclass[10pt,a4paper]{article}
\usepackage[T1]{fontenc}
\usepackage{lmodern}
\usepackage[margin=25mm]{geometry}
\usepackage{amsmath,amssymb,amsthm}
\usepackage[colorlinks=true,urlcolor=blue]{hyperref}
\usepackage{xurl}
\newcommand{\code}[1]{\nolinkurl{#1}}
\newtheorem{proposition}{Proposition}
\newcommand{\Pow}{\mathcal P}
\title{Clopen subsets are not fixed points of complement:\\
a counterexample to conjecture 00000006238}
\author{asfigh665}
\date{October 15, 2026}
\begin{document}
\maketitle

\begin{abstract}
The conjecture asserts that the fixed points of the complement involution
are precisely the clopen elements. In the complete Boolean lattice of
subsets of a discrete two-point space, a singleton is both closed and open,
whereas its complement is the other singleton. The asserted fixed-point
characterization is therefore false, even for a nonempty proper clopen
subset. A standalone Lean 4.24.0 project verifies the lattice completeness,
the topology axioms, the complement operation, and the negation of the
fixed-point assertion.
\end{abstract}

\section{The clause being disproved}
The English statement of conjecture 00000006238 includes the clause
\begin{quote}
``the square of the involution is the identity, with fixed points the
clopen elements.''
\end{quote}
The preceding clause identifies that involution as complement. The Chinese
version likewise states that the involution has the clopen elements as its
fixed points. With $c(A)=X\setminus A$, the disputed characterization is
\begin{equation}\label{eq:claim}
 c(A)=A\quad\Longleftrightarrow\quad
 A\text{ is both closed and open}.
\end{equation}
The asserted identity $c^2=\mathrm{id}$ is valid and is not the point of
failure. A fixed point of $c$ means $c(A)=A$, rather than $c^2(A)=A$.

We use the standard concrete instance of the stated closed--open duality
of subobjects: the powerset lattice of a topological space, with the
usual set complement. In the discrete example below every subset is both
open and closed, so both the open and the closed subobject lattices are
the entire powerset lattice. No domain issue for complement arises.
The conjecture imposes no nontriviality, connectedness, or other condition
excluding this finite example.

\section{An admissible complete lattice and topology}
Let $X=\{0,1\}$ and give it the discrete topology $\tau=\Pow(X)$.
The subobject lattice is
\[
 L=\Pow(X)=\{\varnothing,\{0\},\{1\},X\},
 \qquad A\leq B\ \Longleftrightarrow\ A\subseteq B.
\]
Its meet is intersection, its join is union, and its bottom and top are
$\varnothing$ and $X$. For every family $\mathcal F\subseteq L$, including
the empty family, its least upper bound and greatest lower bound are
\[
 \bigvee\mathcal F=\bigcup_{A\in\mathcal F} A,
 \qquad
 \bigwedge\mathcal F=\bigcap_{A\in\mathcal F} A.
\]
The empty intersection is understood to be $X$. Thus $L$ is complete.
It is a Boolean lattice with the usual complement $c$, and in particular
\[
 A\cap c(A)=\varnothing,\qquad A\cup c(A)=X,
 \qquad c(c(A))=A.
\]
Complement reverses inclusion. Every member of $\tau$ is open by
definition; its complement is also in $\tau$, so every member is closed
as well. The empty set and $X$ are open, and arbitrary unions and finite
intersections of members of $\tau$ remain in $\tau$. These observations
verify the topology axioms directly.

\section{The counterexample}
Take $A=\{0\}$. It is nonempty and proper, and it is clopen in the
discrete topology. Nevertheless
\[
 c(A)=\{1\}\ne\{0\}=A.
\]
Consequently the right-hand side of \eqref{eq:claim} holds and the
left-hand side fails. This disproves that clause and hence the conjunction
asserted in conjecture 00000006238. No claim about the additional,
unspecified ``multiplicative deficit'' is needed.

\begin{proposition}
For every nonempty set $Y$, the complement map on $\Pow(Y)$ has no fixed
point.
\end{proposition}
\begin{proof}
Suppose $B=Y\setminus B$ and choose $y\in Y$. If $y\in B$, the equality
implies $y\notin B$, a contradiction. It follows that $y\notin B$.
But then $y\in Y\setminus B=B$, again a contradiction.
\end{proof}

In particular, clopen means that complement preserves membership in both
the open and closed classes. It does not mean that complement fixes the
individual subset. The four-element lattice already exhibits this
distinction without any infinitary or geometric complications.

\section{Formal verification and semantic correspondence}
The accompanying project uses only Lean's standard library and is pinned
to Lean 4.24.0. It defines a subset of a type $X$ as a predicate
$X\to\mathrm{Prop}$, with inclusion, intersection, union, complement,
arbitrary suprema, and arbitrary infima given by their usual membership
conditions. The theorems \code{supremum_is_lub} and
\code{infimum_is_glb} establish the complete-lattice bounds for
arbitrary families. Theorems for complement establish its two Boolean
identities, reversal of inclusion, and the involution identity.

The structure \texttt{Topology} records that the empty set and the whole
space are open, and that binary intersections and arbitrary unions of
opens are open. The constructed
\texttt{discreteTwo} satisfies all of them. Closed means that the
complement is open, and clopen means open and closed. The two points are
represented by Lean's Boolean values \texttt{false} and \texttt{true};
\texttt{singleton} is the predicate $x=\texttt{false}$.
The source proves that this subset is nonempty, proper, clopen, and not
fixed by complement, and that its complement is exactly
$x=\texttt{true}$. The general theorem \texttt{no\_fixed\_subset}
formalizes the proposition above.

The central result \code{not_discrete_fixed_point_claim}
is the negation of \eqref{eq:claim} restricted to all subsets of this
explicit admissible space. The theorem
\code{not_universal_fixed_point_claim} also refutes the
corresponding statement quantified over topological spaces. Finally,
\code{refutes_every_claim_covering_discrete_case} records the
logical consequence for any larger claim implying the disputed clause
on this example. The project does not assign a new definition to the
conjecture's unspecified deficit term; refuting one explicit required
conjunct suffices.

To reproduce the verification, run in the \texttt{lean4} directory:
\begin{verbatim}
lake build
lake exe verify
lake env lean Conjecture6238.lean
\end{verbatim}
The last command reports the axiom dependencies of the principal results.
The source contains no proof holes, native decision procedure, or added
axioms. All proofs are checked by the Lean kernel. Classical logic is
used for the general double-complement identity; the refutation itself
needs no additional axioms.

\end{document}
